% Avsnitt 10.1, ur lectures/L10/appendix/intro_framing.md.

\section{Ramar och serialisering}
\label{c10:sec:framing}

\Chapref{c1:ch} till \ref{c8:ch} slutade med tillståndsmaskiner som styrde ett enda kort: en
blinkande lysdiod, en UART på en FPGA. Projektet utgår från ett annat problem: att få oberoende
hårdvarudelar att utbyta information tillförlitligt över en delad ledning, där ingen av sidorna kan
se in i den andra.

Det här är centralt i inbyggda system eftersom kommunikationen sker som råa byteströmmar, RAM och
flash är begränsade, och robusthet är kritisk. Noder kommunicerar oftast över seriella gränssnitt
som UART, SPI, I2C eller RS-485. Data skickas som en följd av bytes, och att tolka den byteströmmen
kräver en struktur. Den strukturen kallas en \textbf{ram}.

En ram är en strukturerad representation av data som:
\begin{itemize}
  \item Innehåller tydligt definierade fält.
  \item Har en definierad längd.
  \item Kan valideras (till exempel via en checksumma).
  \item Kan serialiseras till en bytearray (\code{std::uint8_t[]}).
  \item Kan deserialiseras från en bytearray.
\end{itemize}

Betrakta nod A som skickar värdet \code{0x2A} till nod B över en delad ledning, utan att något annat
är överenskommet i förväg:
\begin{itemize}
  \item Hur vet B när ett nytt värde börjar i en ström av bitar som ser likadan ut mitt i ett värde
    som mellan två värden?
  \item Hur vet B hur långt (hur många bytes) värdet är?
  \item Hur vet B att bitarna den samplade inte förvanskades av elektriskt brus?
  \item Vad händer om två noder sänder samtidigt?
\end{itemize}

Varje riktigt protokoll besvarar de här med någon kombination av följande, var och en namngiven här
efter det fält den här bokens ram bär den i:
\begin{itemize}
  \item \textbf{Ramning (\code{SOF})}: en igenkännbar startmarkör (och ofta slutmarkör), så att en
    mottagare kan avgöra var ett meddelande börjar och var det slutar.
  \item \textbf{Ett längdfält (\code{LEN})}: så att mottagaren vet hur mycket nyttolast den ska
    vänta sig.
  \item \textbf{En checksumma (\code{CHK})}: extra bitar, beräknade ur nyttolasten, som låter
    mottagaren upptäcka (inte nödvändigtvis rätta) många sorters överföringsfel.
  \item \textbf{En regel för mediaåtkomst}: vad som händer när fler än en nod vill sända samtidigt.
\end{itemize}

Ett \code{TYPE}-fält sällar sig till dem för att säga vad meddelandet betyder, så att mottagaren kan
agera utifrån det.

\subsection*{En liten men realistisk ram}
För att göra det konkret: konstruera en ram för att skicka korta meddelanden mellan noder på en
delad, byteorienterad länk (UART, RS-485, \dots):

\lecfig[0.68]{lectures/L10/appendix/images/example_frame_format.png}{Ramformatet som resten av det
här avsnittet bygger, fält för fält.}{c10:fig:frame}

Tre regler spelar roll så snart man börjar räkna ut de här fälten för hand:
\begin{itemize}
  \item \textbf{Byteordning}: varje fält bredare än en byte (SOF, SEQ, CHK) är big-endian, mest
    signifikanta byte först.
  \item \textbf{LEN} räknar bara DATA-fältet, aldrig huvudet och aldrig CHK.
  \item \textbf{CHK} är en 16-bitars summa av varje byte från SOF till och med den sista
    nyttolastbyten, som slår runt vid spill. Den beräknas över den \emph{serialiserade}
    byteströmmen snarare än över fältvärdena som heltal, så SOF bidrar med
    \code{0xA5 + 0xF7}, inte \code{0xA5F7} behandlat som ett enda 16-bitarstal.
\end{itemize}

\paragraph{Varför just DST, SRC och SEQ}
SOF, LEN, TYPE och CHK besvarar redan de fyra frågorna ovan. De här tre finns eftersom en delad länk
med fler än två deltagare väcker nya:
\begin{itemize}
  \item \textbf{DST och SRC}: vem ramen är till och vem den är från, så att en nod kan ignorera
    trafik adresserad någon annanstans och ett svar kan hitta tillbaka. CAN besvarar det här helt
    annorlunda, som \secref{c10:sec:broadcast} visar.
  \item \textbf{SEQ}: matchar ett svar tillbaka mot \emph{just den} förfrågan, inte någon annan som
    är i luften.
\end{itemize}

\subsection*{Att välja SOF}
\begin{itemize}
  \item En enda synkbyte är lätt att träffa av misstag: vilken nyttolastbyte som helst som råkar
    vara lika med den ser ut som början på en ram.
  \item Två bytes är därför vanligt, och det här protokollet använder \code{0xA5F7}.
  \item SOF är det \emph{första} en parser kontrollerar när den läser en ram. Stämmer den inte
    kastas datan, och parsern fortsätter söka efter nästa SOF.
  \item Lägg märke till vad det ger och inte ger: det gör en falsk synkronisering mycket mindre
    sannolik, men inte omöjlig, eftersom nyttolastbytes fortfarande kan råka stava \code{A5 F7}. Den
    kvarvarande risken är varför LEN och CHK ändå måste kontrolleras efteråt.
\end{itemize}

\subsection*{Genomräknat exempel: PING och PONG}
Anta att nod \code{0x01} pingar nod \code{0x02}, med sekvensnummer \code{0x0020} och utan nyttolast:
\begin{itemize}
  \item \code{SOF = 0xA5F7}, \code{LEN = 0x00}, \code{TYPE = 0x00} (Ping), \code{DST = 0x02},
    \code{SRC = 0x01}, \code{SEQ = 0x0020}.
  \item \code{CHK = 0xA5 + 0xF7 + 0x00 + 0x00 + 0x02 + 0x01 + 0x00 + 0x20 = 0x01BF}.
\end{itemize}

\begin{textdrawing}
SOF      LEN   TYPE   DST   SRC   SEQ      DATA   CHK
A5 F7    00    00     02    01    00 20    --     01 BF
\end{textdrawing}

Nod \code{0x02} svarar med en PONG:
\begin{itemize}
  \item \code{SOF = 0xA5F7}, \code{LEN = 0x00}, \code{TYPE = 0x01} (Pong), \code{DST = 0x01},
    \code{SRC = 0x02}, \code{SEQ = 0x0020}.
  \item \code{CHK = 0xA5 + 0xF7 + 0x00 + 0x01 + 0x01 + 0x02 + 0x00 + 0x20 = 0x01C0}.
  \item Lägg märke till att \code{DST} och \code{SRC} är omkastade, \code{SEQ} är oförändrad, och
    bara \code{TYPE} och de två adresserna skiljer sig från PING:en: identisk aritmetik med andra
    indata.
\end{itemize}

Det täcker tre av de fyra inledande frågorna, plus de extra som en delad länk med flera noder
väcker. Den fjärde står fortfarande öppen: vad händer om två noder sänder samtidigt? CAN besvarar
den annorlunda, och det är ämnet för \secref{c10:sec:can}.

\subsection*{Ramtyper}
TYPE är en byte på ledningen, så en ramtyp färdas som ett teckenlöst heltal. Varje typ tilldelas
sitt eget id enligt konvention: \code{Ping = 0}, \code{Pong = 1}, \code{StatusReq = 2},
\code{StatusResp = 3}. Ytterligare typer som protokollet får fortsätter helt enkelt sekvensen.

Två följder värda att notera innan någon kod skrivs:
\begin{itemize}
  \item Den mottagna byten är obetrodd indata. Varje värde mottagaren inte känner igen måste
    förkastas i stället för att agera på, precis som en felaktig SOF eller en felaktig checksumma.
  \item Att reservera ett ID direkt ovanför den sista riktiga typen, ett ”okänt”-värde, gör den
    kontrollen till en enda jämförelse. Utan det behöver valideringen en lista över giltiga ID:n som
    måste hållas i takt varje gång en typ läggs till.
\end{itemize}

\subsection*{En enkel 16-bitars checksumma}
För den här bokens syfte behöver en checksumma bara vara enkel att implementera, enkel att testa,
och tillräckligt bra för att fånga många slumpmässiga bitfel. En 16-bitars summa uppfyller alla
tre: börja på \code{0}, och för varje byte i intervallet, \code{sum += byte}. Resultatet är
checksumman, modulo $2^{16}$.

\begin{cppcode}
/**
 * @brief Compute a simple 16-bit checksum (sum of bytes modulo 2^16).
 *
 * @param[in] data Pointer to the bytes to checksum.
 * @param[in] dataLen Number of bytes to checksum.
 *
 * @return The 16-bit checksum.
 */
std::uint16_t checksum16(const std::uint8_t* data, const std::size_t dataLen) noexcept
{
    if ((nullptr == data) || (0U == dataLen)) { return 0U; }
    std::uint16_t checksum{};
    for (std::size_t i{}; i < dataLen; ++i) { checksum += data[i]; }
    return checksum;
}
\end{cppcode}

\begin{warning}
\textbf{En checksummas begränsningar, värda att vara tydlig med.} En checksumma upptäcker bara fel
som ändrar den aritmetiska summan. En enda vänd bit ändrar alltid summan, så enbitsfel fångas
alltid. De riktiga svagheterna ligger någon annanstans: två hela bytes som byter plats lämnar summan
oförändrad, och \emph{vissa} kombinationer av flera vända bitar tar ut varandra och slinker igenom
oupptäckta. CRC:er fångar en mycket bredare, matematiskt välförstådd klass av fel, och används därför
av CAN i stället för en enkel summa. De introduceras i \secref{c10:sec:frameformat} och
implementeras i \chapref{c13:ch}.
\end{warning}

Det här är inte lika starkt som en CRC, men det räcker för att öva ramning, parsning och robust
felhantering, vilket är vad den här introduktionen är till för.

\subsection*{Serialisering och deserialisering}
\textbf{Serialisering} omvandlar strukturerad data till en följd av bytes (\code{std::uint8_t[]}) så
att den kan skickas över ett gränssnitt. \textbf{Deserialisering} tolkar en mottagen bytearray och
plockar ut den strukturerade datan igen: fälten valideras först (SOF, längd, typ, checksumma), och
först när varje kontroll passerat plockas datan ut.

Att validera före utplockning spelar roll eftersom UART och RS-485 levererar bytes som en ström, så
en mottagare rutinmässigt ser skräp mellan ramar, tappade bytes och förvanskade ramar.

I inbyggda system görs det här nästan alltid med C-arrayer (\code{std::uint8_t buffer[Size]}),
pekare till de arrayerna vid anropsgränser (\code{std::uint8_t*}), och en explicit längd
(\code{std::size_t}).

Båda riktningarna behöver samma hantering av byteordning på varje flerbytesfält, vilket är
repetitivt och lätt att få subtilt fel. Två funktionsmallar tar bort upprepningen. Båda begränsar
\code{T} med \code{std::is_unsigned}, så var de än hamnar behöver de \code{<type_traits>}:

\begin{cppcode}
#include <type_traits>

/**
 * @brief Write a value to a buffer, most significant byte first.
 *
 *        The caller is responsible for the buffer being large enough.
 *
 * @tparam T The value type. Must be unsigned.
 *
 * @param[out] buf The buffer to write to.
 * @param[in] offset Value offset, i.e. the index at which to write the value.
 * @param[in] value  The value to write.
 */
template <typename T>
void writeToBuf(std::uint8_t* buf, const std::size_t offset, const T value) noexcept
{
    static_assert(std::is_unsigned<T>::value, "Failed to write to buffer: T must be unsigned!");
    constexpr std::size_t len{sizeof(T)};
    constexpr std::size_t msb{len - 1U};
    constexpr std::size_t bitsPerByte{8U};

    for (std::size_t i{}; i < len; ++i)
    {
        const std::size_t shift{bitsPerByte * (msb - i)};
        buf[offset + i] = static_cast<std::uint8_t>((value >> shift));
    }
}

/**
 * @brief Read a value from a buffer, most significant byte first.
 *
 * @tparam T The value type. Must be unsigned.
 *
 * @param[in] buf Buffer to read from.
 * @param[in] offset Value offset in the buffer.
 *
 * @return The retrieved value.
 */
template <typename T = std::uint8_t>
[[nodiscard]] T readFromBuf(const std::uint8_t* buf, const std::size_t offset) noexcept
{
    static_assert(std::is_unsigned<T>::value, "Failed to read from buffer: T must be unsigned!");
    constexpr std::size_t len{sizeof(T)};
    constexpr std::size_t msb{len - 1U};
    constexpr std::size_t bitsPerByte{8U};
    T value{};

    for (std::size_t i{}; i < len; ++i)
    {
        const std::size_t shift{bitsPerByte * (msb - i)};
        value |= static_cast<T>(buf[offset + i]) << shift;
    }
    return value;
}
\end{cppcode}

\begin{warning}
\textbf{En fälla värd att namnge, eftersom kompilatorn inte fångar den.} Antalet bytes som skrivs är
\code{sizeof(T)}, härlett ur \emph{argumentet}, inte ur fältet. Att skicka in en \code{std::size_t}
där fältet är en byte skriver \textbf{åtta} bytes och skriver i tysthet över fälten som följer:

\begin{cppcode}
constexpr std::size_t payloadLen{0U};
writeToBuf(buf, FrameOffset::Len, payloadLen);                            // Writes 8 bytes (error).
writeToBuf(buf, FrameOffset::Len, static_cast<std::uint8_t>(payloadLen)); // Writes 1 byte.
\end{cppcode}

Casta alltid till fältets egen bredd på anropsstället.
\end{warning}

\subsection*{Från papper till kod}
Precis den här ramen görs körbar i C++ under \file{lectures/L10/appendix/code/}: en \code{Frame}-\!
struct som serialiserar sig själv till en bytebuffert, deserialiserar sig själv tillbaka, och
förkastar allt med felaktig SOF, längd, typ eller checksumma. Att serialisera PING:en ovan
producerar exakt de bytes som räknades ut ovan, \code{A5 F7 00 00 02 01 00 20 01 BF}.
\file{comm/def.hpp} tillhandahåller redan fältoffseten och fältlängderna, så du behöver aldrig räkna
bytepositioner för hand.

Att göra typ-ID:na ovan till en C++-uppräkning, deklarera structen och implementera båda metoderna
själv är Övning~\ref{c10:ex:cpp}. En uppsättning enhetstester följer med koden, i
\file{test/frame_test.cpp}: de kontrollerar båda metoderna mot ramarna som räknades ut på papper,
byte för byte. Du skriver dem inte; \code{make} kör dem, och din uppgift är att få dem att passera.

\subsection*{Vad som kommer härnäst}
En fråga från början av det här avsnittet står fortfarande öppen: \textbf{vad händer om två noder
sänder samtidigt?}

Ingenting i det här ramformatet besvarar den. Ett synkord kan inte hjälpa, eftersom båda sändarna
skickar ett; en checksumma kan inte hjälpa, eftersom den bara rapporterar skadan i efterhand. Att
besvara den kräver ett annat slags protokoll, ett där bussens \emph{elektriska} beteende avgör vem
som fortsätter och vem som faller ifrån, innan någon data går förlorad. Det protokollet är CAN.
